\section{Conclusion}
\label{sec:conclusion}

Reviews of the AD5940-based low-cost potentiostat literature, including the
most recent one, evaluate control software functionally and leave its
architecture unexamined. We audited one such instrument's firmware, found
concrete, fixable problems -- triplicated measurement-configuration code,
two entirely dead parallel subsystems, and a silent status-indicator bug --
and fixed them with a scoped, individually verifiable refactor. Reusing the
dummy-cell validation protocol from a companion metrological study, we
measured, rather than assumed, that the refactor left the instrument's
CV, CA, SWV, and DPV output effectively unchanged (mean relative CV
difference $0.070\,\%$), and corroborated the result on an independent real
load. Because this instrument's analog front end is commercial,
characterized silicon, the residual reliability risk in AD5940-based
potentiostats is concentrated in firmware structure rather than analog
design -- and that risk, this work shows, can be reduced with ordinary
refactoring technique at effectively no cost to measurement correctness.
We see this as a first, scoped demonstration of a broader point: for this
growing instrument class, software architecture deserves the same
deliberate scrutiny that these papers already give to transimpedance gain
tables and calibration protocols, and doing so does not have to come at the
expense of the numbers the instrument was built to report.
